
\documentclass[11pt,a4paper]{article}
\usepackage[margin=0.85in]{geometry}
\usepackage[T1]{fontenc}
\usepackage[utf8]{inputenc}
\usepackage{lmodern}
\usepackage{booktabs}
\usepackage{longtable}
\usepackage{array}
\usepackage{amsmath}
\usepackage{xcolor}
\usepackage{listings}
\usepackage[hidelinks,breaklinks=true]{hyperref}

\definecolor{codebg}{RGB}{248,248,248}
\definecolor{codeframe}{RGB}{220,220,220}
\newcommand{\codepath}[1]{\path{#1}}
\newcommand{\na}{\textsc{NA}}

\lstdefinestyle{cmdstyle}{
  basicstyle=\ttfamily\small,
  breaklines=true,
  columns=fullflexible,
  backgroundcolor=\color{codebg},
  frame=single,
  rulecolor=\color{codeframe},
  keepspaces=true,
  showstringspaces=false
}

\title{LTE VANET 150-Second Comparison Report: ECC, PBC, and PBC with Federated Learning}
\author{VANET ns-3 Security Research Workspace}
\date{May 3, 2026}

\begin{document}
\maketitle

\begin{abstract}
This report compares three LTE VANET simulation cases over the same 150-second
run window: normal ECC authentication, PBC authentication, and PBC authentication
with federated learning (FL). The goal is to quantify how the stronger PBC
security model and the added FL workflow affect latency, energy, average power,
communication overhead, and authentication behavior. The results show that ECC
has the lowest total energy and lowest cryptographic latency. PBC significantly
increases security energy and verification latency because pairing-based
aggregate verification is expensive. Adding FL on top of PBC adds a small direct
FL-only communication/computation cost, while the total PBC+FL run shows higher
security energy than PBC alone. In the completed FL round, 10 out of 10
participant updates were authenticated successfully.
\end{abstract}

\tableofcontents
\newpage

\section{Experiment Scope}

The three cases are:

\begin{longtable}{p{0.22\linewidth}p{0.68\linewidth}}
\toprule
\textbf{Case} & \textbf{Description} \\
\midrule
LTE ECC & Baseline LTE VANET using ECC/ECDSA authentication. Vehicles receive ECC credentials, sign messages, and receivers verify ECC signatures. \\
LTE PBC & LTE VANET using PBC security. Vehicles receive PBC pseudonym/key material, sign relay-mode packets, and the relay performs aggregate verification. \\
LTE PBC + FL & LTE PBC security with federated learning enabled. Vehicles receive a global model, compute local updates, sign FL updates, RSUs verify/aggregate updates, and the BS performs global aggregation. \\
\bottomrule
\end{longtable}

All three cases were run for the same requested simulation duration:
\[
T_{sim}=150\text{ s}.
\]
The final periodic CSV row is at simulation time 149 s, which is expected because
the logger writes rows once per simulated second before the simulator ends.

\section{Commands Used}

\begin{lstlisting}[style=cmdstyle,caption={Workspace and build command.}]
cd /home/sahib/vanet_ns3_security_research/simulation_workspace/ns-allinone-3.44/ns-3.44
mkdir -p results/lte_3case_150
./ns3 build vanet-security-lte vanet-security-nr
\end{lstlisting}

\begin{lstlisting}[style=cmdstyle,caption={LTE ECC run.}]
LSAN_OPTIONS=detect_leaks=0 ./ns3 run "vanet-security-lte --securityMode=ecc --networkLabel=5g_lte_ecc --effectiveRateMbps=100 --simTime=150 --progressLog=true --outCsv=results/lte_3case_150/lte_ecc_power.csv --stageCsv=results/lte_3case_150/lte_ecc_stages.csv"
\end{lstlisting}

\begin{lstlisting}[style=cmdstyle,caption={LTE PBC run.}]
LSAN_OPTIONS=detect_leaks=0 ./ns3 run "vanet-security-lte --securityMode=pbc --networkLabel=5g_lte_pbc --effectiveRateMbps=100 --simTime=150 --progressLog=true --outCsv=results/lte_3case_150/lte_pbc_power.csv --stageCsv=results/lte_3case_150/lte_pbc_stages.csv"
\end{lstlisting}

\begin{lstlisting}[style=cmdstyle,caption={LTE PBC + federated learning run.}]
LSAN_OPTIONS=detect_leaks=0 ./ns3 run "vanet-security-lte --securityMode=pbc --networkLabel=5g_lte_pbc_fl --effectiveRateMbps=100 --enableFl=true --flSecurityMode=pbc --flRounds=3 --flParticipants=10 --flModelDim=16 --simTime=150 --progressLog=true --outCsv=results/lte_3case_150/lte_pbc_fl_power.csv --stageCsv=results/lte_3case_150/lte_pbc_fl_stages.csv --flCsv=results/lte_3case_150/lte_pbc_fl_rounds.csv --flStageCsv=results/lte_3case_150/lte_pbc_fl_stage_summary.csv"
\end{lstlisting}

\section{Main Power and Latency Results}

\begin{longtable}{p{0.18\linewidth}rrrrrr}
\toprule
\textbf{Case} & \textbf{Failures} & \textbf{Security J} & \textbf{Comm J} & \textbf{Total J} & \textbf{Avg Power W} & \textbf{V2V ms} \\
\midrule
LTE ECC & 0 & 59.437 & 19.511 & 78.947 & 0.530 & 849.767 \\
LTE PBC & 156 & 383.359 & 5.928 & 389.288 & 2.613 & 2204.380 \\
LTE PBC + FL & 156 & 413.805 & 5.935 & 419.741 & 2.817 & 2197.090 \\
\bottomrule
\end{longtable}

\begin{longtable}{p{0.20\linewidth}rrrrrr}
\toprule
\textbf{Case} & \textbf{V2I Up ms} & \textbf{V2I Down ms} & \textbf{V2I RTT ms} & \textbf{Reg ms} & \textbf{Sign ms} & \textbf{Verify ms} \\
\midrule
LTE ECC & 33.855 & 812.608 & 846.599 & 280.780 & 1.551 & 1.594 \\
LTE PBC & 32.938 & 524.476 & 556.106 & 495.800 & 27.458 & 111.197 \\
LTE PBC + FL & 32.674 & 522.515 & 554.148 & 495.800 & 25.816 & 120.773 \\
\bottomrule
\end{longtable}

\section{Stage-Level Summary}

\begin{longtable}{p{0.20\linewidth}rrrrrr}
\toprule
\textbf{Case} & \textbf{Security J} & \textbf{Security ms} & \textbf{Comm J} & \textbf{Comm ms} & \textbf{FL-only J} & \textbf{FL events} \\
\midrule
LTE ECC & 59.789 & 23384.777 & 19.636 & 20944.672 & 0.000000 & 0 \\
LTE PBC & 383.359 & 26319.416 & 5.929 & 6223.433 & 0.000000 & 0 \\
LTE PBC + FL & 413.805 & 28778.050 & 5.936 & 6226.308 & 0.007113 & 134 \\
\bottomrule
\end{longtable}

\section{Increment Analysis}

\subsection{PBC Compared with ECC}

\begin{longtable}{p{0.42\linewidth}rr}
\toprule
\textbf{Metric} & \textbf{Absolute change} & \textbf{Percent change} \\
\midrule
Security energy & 323.922 J & 544.99\% \\
Communication energy & -13.582 J & -69.61\% \\
Total energy & 310.341 J & 393.10\% \\
Average total power & 2.083 W & 393.10\% \\
Average V2V delay & 1354.613 ms & 159.41\% \\
Average V2I RTT & -290.493 ms & -34.31\% \\
Registration delay & 215.020 ms & 76.58\% \\
Signing latency & 25.907 ms & 1670.51\% \\
Verification latency & 109.603 ms & 6877.64\% \\
\bottomrule
\end{longtable}

\textbf{Interpretation.} PBC is much heavier than ECC in this run. Total energy
increased by 310.341 J, and average total
power increased by 2.083 W. The main reason is
cryptographic cost: average signing latency rose from 1.551 ms to
27.458 ms, and average verification latency rose from
1.594 ms to 111.197 ms. The PBC
registration path is also heavier, with registration delay increasing from
280.780 ms to 495.800 ms.

\subsection{PBC + FL Compared with PBC}

\begin{longtable}{p{0.42\linewidth}rr}
\toprule
\textbf{Metric} & \textbf{Absolute change} & \textbf{Percent change} \\
\midrule
Security energy & 30.446 J & 7.94\% \\
Communication energy & 0.006670 J & 0.11\% \\
Total energy & 30.453 J & 7.82\% \\
Average total power & 0.204 W & 7.82\% \\
Average V2V delay & -7.290 ms & -0.33\% \\
Average V2I RTT & -1.958 ms & -0.35\% \\
Signing latency & -1.642 ms & -5.98\% \\
Aggregate latency & 0.017 ms & 9.13\% \\
Verification latency & 9.576 ms & 8.61\% \\
Partial-key latency & 2.191 ms & 41.07\% \\
\bottomrule
\end{longtable}

\textbf{Interpretation.} Adding FL on top of PBC increased total energy by
30.453 J, or 7.82\%.
The direct communication-energy increase was very small: only
0.006670 J. Stage-level FL-only overhead was
0.007113 J across 134
FL-specific events, with 2.985927 ms of FL-only latency.
This means the FL traffic itself is visible but small compared with the PBC
verification workload.

\section{Federated Learning Authentication Result}

The FL-enabled case completed one global FL round. The round-level output was:

\begin{longtable}{p{0.42\linewidth}p{0.40\linewidth}}
\toprule
\textbf{FL metric} & \textbf{Value} \\
\midrule
Round completed & 0 \\
Completion time & 10.402010 s \\
Selected vehicles & 10 \\
RSU count & 2 \\
Updates sent & 10 \\
Updates verified & 10 \\
Updates rejected & 0 \\
Authentication success rate & 100\% \\
Model download bytes & 1632 \\
Update upload bytes & 1440 \\
Average FL update signing latency & 22.221816 ms \\
Average FL update verification latency & 115.972422 ms \\
\bottomrule
\end{longtable}

\section{Vehicle Credential and Authentication Counts}

The following table separates three ideas that are easy to confuse: total
vehicles that received credentials, vehicles that generated PBC full-key
material, and vehicles selected for the completed FL training round.

\begin{longtable}{p{0.25\linewidth}p{0.28\linewidth}p{0.35\linewidth}}
\toprule
\textbf{Case} & \textbf{Observed count} & \textbf{Meaning} \\
\midrule
LTE ECC & \codepath{vehicle:register_resp}=50 & 50 vehicles received ECC registration responses. This means the baseline ECC credential-delivery path completed for 50 vehicles. \\
LTE PBC & \codepath{vehicle:register_resp}=50 & 50 vehicles received PBC-mode registration responses through the vehicle registration path. \\
LTE PBC & \codepath{pbc_full_key}=50 & 50 vehicles generated PBC full-key material after receiving PBC credentials. \\
LTE PBC + FL & \codepath{vehicle:register_resp}=50 & 50 normal VANET vehicles received PBC-mode registration responses. \\
LTE PBC + FL & \codepath{pbc_full_key}=60 & 60 PBC full-key generation events were recorded. This is not 60 unique VANET cars; it includes the normal 50 VANET vehicle credential events plus the additional FL participant/security app credential events. \\
LTE PBC + FL & \codepath{selected_vehicles}=10 & 10 vehicles were selected to participate in the completed FL round. \\
LTE PBC + FL & \codepath{updates_verified}=10 & All 10 submitted FL updates were authenticated successfully. \\
LTE PBC + FL & \codepath{updates_rejected}=0 & No FL update was rejected. \\
\bottomrule
\end{longtable}

\textbf{Result.} The completed FL round authenticated 10 out of 10 participating
vehicle updates:
\[
\text{FL authentication success rate}=\frac{10}{10}\times100=100\%.
\]

\textbf{Important distinction.} The FL update authentication result is clean:
10 out of 10 FL updates were verified and 0 were rejected. However, the periodic
VANET-wide metric \codepath{verification_failures} is 156 for both PBC and
PBC+FL. Because the value is identical in PBC and PBC+FL, FL did not add extra
VANET-wide verification failures in this run. These 156 failures belong to the
broader PBC VANET relay/verification workload, not to rejected FL updates.

\section{Meaning of the 156 Verification Failures}

The value \codepath{verification_failures=156} does \emph{not} mean that 156
vehicles failed authentication. It is a cumulative packet/event counter. In the
implementation, \codepath{VanetStats::RecordVerificationFailure()} increments
one global counter whenever a security verification guard fails during packet
processing. In PBC relay mode, this can happen when a packet has invalid or
missing PBC authentication fields, a PID timestamp is outside the accepted
freshness window, PID encoding fails, an aggregate batch fails pairing
verification, or a vehicle receives a PBC packet from a non-trusted relay path.

For the PBC relay, one failed aggregate batch can add more than one failure,
because the relay records one failure for each packet entry in the rejected
batch. Therefore, the value 156 should be read as:

\begin{quote}
156 packet-level or batch-entry verification failure events occurred during the
150-second PBC VANET message-propagation workload.
\end{quote}

It should \textbf{not} be read as:

\begin{quote}
156 vehicles failed to register, or 156 FL updates were rejected.
\end{quote}

The evidence is:

\begin{itemize}
  \item LTE PBC and LTE PBC+FL both have \codepath{verification_failures}=156.
  \item The FL round has \codepath{updates_sent}=10, \codepath{updates_verified}=10, and \codepath{updates_rejected}=0.
  \item Therefore, the FL layer did not add new rejected updates in this run.
  \item The counter is coming from the broader PBC VANET relay/verification workload, not from the completed FL update-authentication path.
\end{itemize}

\section{Main Findings}

\begin{itemize}
  \item \textbf{ECC is the lowest-power baseline.} LTE ECC consumed 78.947 J total energy and 0.530 W average total power.
  \item \textbf{PBC costs more than ECC.} LTE PBC consumed 389.288 J, which is 310.341 J higher than ECC.
  \item \textbf{PBC verification is the dominant security bottleneck.} PBC average verification latency was 111.197 ms, compared with 1.594 ms for ECC.
  \item \textbf{Adding FL to PBC increased total energy moderately.} LTE PBC+FL consumed 419.741 J, which is 30.453 J higher than PBC alone.
  \item \textbf{Direct FL communication overhead is small.} FL-only stage energy is 0.007113 J and FL-only latency is 2.985927 ms.
  \item \textbf{FL authentication succeeded.} In the completed FL round, 10 out of 10 vehicle updates were verified successfully.
\end{itemize}

\section{Conclusion}

The 150-second LTE comparison shows a clear progression of cost. ECC is the most
lightweight authentication mode. PBC introduces substantially higher security
energy and latency, especially during aggregate verification, but it provides the
pairing-based authenticated relay model required by the project. Adding FL on
top of PBC adds measurable but comparatively small direct FL communication and
computation overhead. The total PBC+FL energy is higher than PBC alone, but the
FL-specific stage overhead is small in the current ns-3 implementation because
local training is represented by a lightweight synthetic vector update. For a
future experiment that highlights training delay more strongly, the simulator
should add a configurable training-time parameter or import measured training
latency from the Python FL baseline.

\section{Files Used}

\begin{itemize}
  \item \codepath{simulation_workspace/ns-allinone-3.44/ns-3.44/results/lte_3case_150/lte_3case_power_latency_summary.csv}
  \item \codepath{simulation_workspace/ns-allinone-3.44/ns-3.44/results/lte_3case_150/lte_3case_stage_summary.csv}
  \item \codepath{simulation_workspace/ns-allinone-3.44/ns-3.44/results/lte_3case_150/lte_pbc_fl_rounds.csv}
\end{itemize}

\end{document}
